% ---- score2 standard preamble (from: python3 score2_tools.py preamble) ----
\documentclass[10pt]{extarticle}
\usepackage[a4paper,margin=1.5cm]{geometry}
\usepackage[T1]{fontenc}
\usepackage{amsmath,amssymb}
\usepackage{graphicx}
\usepackage[export]{adjustbox}
\usepackage{multicol}
\usepackage{enumitem}
\usepackage{booktabs,array,tabularx}
\setlength{\columnsep}{0.7cm}
\setlength{\parindent}{0pt}
\setlength{\parskip}{3pt}
\pagestyle{empty}
\makeatletter
% figure registry (figdefs.tex is generated by the extract step)
\newcommand{\qfigdef}[4]{\@namedef{qf@file@#1}{#2}\@namedef{qf@w@#1}{#3}\@namedef{qf@wide@#1}{#4}}
\newcommand{\qf@check}[1]{\@ifundefined{qf@file@#1}{\PackageError{score2}{Unknown figure id #1}{}}{}}
\newcommand{\qf@img}[4][]{\includegraphics[width=#2\dimexpr\@nameuse{qf@w@#4}pt\relax,max width=#3,max height=0.45\textheight,#1]{\@nameuse{qf@file@#4}}}
% Full-width questions: finalize wraps every question that contains a wide figure in \qfullbegin...\qfullend,
% so the WHOLE question (stem, figures, options) is set across the page in reading order, then the two columns
% resume. The question's head - its number, stem and first figure (\qfullsplit marks the end of the head) -
% is kept on one page, so a number is never left at a page foot with its figure overleaf; the rest may break
% across pages normally, so no big gaps are left. Without a \qfullsplit, the whole question is kept together
% if it is under 45% of a page. Never write these yourself.
\newif\ifqf@full
\newif\ifqf@boxed
\newsavebox{\qf@box}
\newcommand{\qf@open}{\global\qf@fulltrue\par\global\qf@boxedtrue
  \setbox\qf@box\vbox\bgroup\hsize\textwidth\linewidth\textwidth}
\newcommand{\qf@emit}[1]{\par\egroup\global\qf@boxedfalse
  \ifdim\dimexpr\ht\qf@box+\dp\qf@box\relax<#1\textheight\noindent\usebox{\qf@box}\par\else\unvbox\qf@box\fi}
\newcommand{\qfullsplit}{\ifqf@boxed\qf@emit{0.85}\fi}
\newcommand{\qf@close}{\ifqf@boxed\qf@emit{0.45}\fi\global\qf@fullfalse}
\newcommand{\qfullbegin}{\end{multicols}\qf@open}
\newcommand{\qfullend}{\qf@close\begin{multicols}{2}}
\newcommand{\qfullstart}{\qf@open}
\newcommand{\qfulldone}{\qf@close}
% \qfig[scale]{id}: block figure, centred, directly where it is written (size set by the tool).
% Fallback if a wide figure is not inside a full-width question: break out of the columns in place - no floats.
\newcommand{\qfig}[2][1]{\qf@check{#2}%
  \ifqf@full\qf@blk{#1}{#2}\else
  \ifnum\@nameuse{qf@wide@#2}=1 \expandafter\@firstoftwo\else\expandafter\@secondoftwo\fi
  {\end{multicols}\begin{center}\qf@img{#1}{\textwidth}{#2}\end{center}\begin{multicols}{2}}%
  {\qf@blk{#1}{#2}}\fi}
\newcommand{\qf@blk}[2]{\par\smallskip{\centering\qf@img{#1}{\linewidth}{#2}\par}\smallskip}
% \qfigi[scale]{id}: inline figure (inside option lists / tables / minipages); never breaks columns.
\newcommand{\qfigi}[2][1]{\qf@check{#2}\qf@img[valign=c]{#1}{\linewidth}{#2}}
\makeatother
% structure macros - lint reads these to check completeness and figure placement
\newcommand{\papersec}[1]{\par\bigskip{\centering\Large\bfseries PAPER-#1\par}\medskip}
\newcommand{\partsec}[2]{\par\medskip\noindent{\large\bfseries PART-#1: #2}\par\smallskip}
\newcommand{\sect}[1]{\par\smallskip\noindent\textbf{#1}\par}
\newcommand{\qn}[1]{\par\smallskip\noindent\textbf{#1.}~}
\newcommand{\paraq}[2]{\par\smallskip\noindent\textbf{Paragraph for Questions #1 and #2}\par}
\newenvironment{qoptions}{\begin{itemize}[leftmargin=1.8em,itemsep=1pt,topsep=2pt,label={}]}{\end{itemize}}
\newcommand{\opt}[2]{\item[(#1)] #2}
% figure registry (generated by score2_tools.py finalize)
\qfigdef{p1_phy_s1_q11_1}{Score-II - Test 05 (22 Mar 2022)-img/p1_phy_s1_q11_1.png}{143.8}{0}
\qfigdef{p1_phy_s1_q12_1}{Score-II - Test 05 (22 Mar 2022)-img/p1_phy_s1_q12_1.png}{89.0}{0}
\qfigdef{p1_chem_s1_q10_A}{Score-II - Test 05 (22 Mar 2022)-img/p1_chem_s1_q10_A.png}{92.0}{0}
\qfigdef{p1_chem_s1_q10_B}{Score-II - Test 05 (22 Mar 2022)-img/p1_chem_s1_q10_B.png}{84.6}{0}
\qfigdef{p1_chem_s1_q10_C}{Score-II - Test 05 (22 Mar 2022)-img/p1_chem_s1_q10_C.png}{91.0}{0}
\qfigdef{p1_chem_s1_q10_D}{Score-II - Test 05 (22 Mar 2022)-img/p1_chem_s1_q10_D.png}{91.0}{0}
\qfigdef{p1_chem_s1_q11_1}{Score-II - Test 05 (22 Mar 2022)-img/p1_chem_s1_q11_1.png}{166.2}{0}
\raggedcolumns
\makeatletter\@beginparpenalty=10000\makeatother
\begin{document}
\begin{center}{\Large\bfseries JEE (Main) Score-II --- Test 05 (22 Mar 2022)}\end{center}
\begin{center}\small\itshape Note: the source file for this test is an incomplete extract (bilingual pages 5, 14, 21--24 of 28 only); the English text of every question on those pages is reproduced below.\end{center}
\begin{multicols}{2}

\partsec{1}{PHYSICS}
\sect{Section-I}

\qn{10} What is the number of $\alpha$ and $\beta$- particle emitted in the following radioactive decay ?\\
${}_{90}X^{200}\rightarrow{}_{80}Y^{168}$
\begin{qoptions}
\opt{A}{6 and 6}
\opt{B}{8 and 8}
\opt{C}{8 and 6}
\opt{D}{6 and 8}
\end{qoptions}

\qn{11} Find output Y :
\qfig{p1_phy_s1_q11_1}
\begin{qoptions}
\opt{A}{AB}
\opt{B}{$\overline{AB}$}
\opt{C}{$\bar AB+A\bar B$}
\opt{D}{$\overline{A+B}$}
\end{qoptions}

\qn{12} Three charges Q, $+$q and $+$q are placed at the verticles of a right -- angled isosceles triangle as shown. The net electrostatic energy of the configuration is zero if Q is equal to :-
\qfig{p1_phy_s1_q12_1}
\begin{qoptions}
\opt{A}{$\dfrac{-q}{1+\sqrt2}$}
\opt{B}{$\dfrac{-2q}{2+\sqrt2}$}
\opt{C}{$-2q$}
\opt{D}{$+q$}
\end{qoptions}

\partsec{2}{CHEMISTRY}
\sect{Section-I}

\qn{10} Among the following the least stable resonance structure is :
\begin{qoptions}
\opt{A}{\qfigi{p1_chem_s1_q10_A}}
\opt{B}{\qfigi{p1_chem_s1_q10_B}}
\opt{C}{\qfigi{p1_chem_s1_q10_C}}
\opt{D}{\qfigi{p1_chem_s1_q10_D}}
\end{qoptions}

\qn{11} In the given transformation, which of the following is the most appropriate reagent ?
\qfig{p1_chem_s1_q11_1}
\begin{qoptions}
\opt{A}{Zn-Hg/HCl}
\opt{B}{Na Liq NH$_3$}
\opt{C}{NaBH$_4$}
\opt{D}{NH$_2$NH$_2$, OH$^-$}
\end{qoptions}

\partsec{3}{MATHEMATICS}
\sect{Section-I}

\qn{6} If $\alpha,\beta,\gamma$ are the real roots of the equation $x^3-3px^2+3qx-1=0$, then the centroid of the triangle whose vertices are $\left(\alpha,\dfrac1\alpha\right)$, $\left(\beta,\dfrac1\beta\right)$ and $\left(\gamma,\dfrac1\gamma\right)$ is :-
\begin{qoptions}
\opt{A}{p, $-$q}
\opt{B}{($-$p, q)}
\opt{C}{(p, q)}
\opt{D}{$\left(\dfrac p2,\dfrac q2\right)$}
\end{qoptions}

\qn{7} $(13)^{507}$ when divided by 9 leaves the remainder :-
\begin{qoptions}
\opt{A}{1}
\opt{B}{4}
\opt{C}{5}
\opt{D}{7}
\end{qoptions}

\qn{8} If $S=1^2+3^2+5^2+\dots+(99)^2$ then the value of $2^2+4^2+6^2+\dots(100)^2$ is :-
\begin{qoptions}
\opt{A}{S + 2550}
\opt{B}{2S}
\opt{C}{4S}
\opt{D}{S + 5050}
\end{qoptions}

\qn{9} The means of five observations is 4 and their variance is 5.2. If three of these observations are 1, 2 and 6, then the other two are :
\begin{qoptions}
\opt{A}{2 and 9}
\opt{B}{3 and 8}
\opt{C}{4 and 7}
\opt{D}{5 and 6}
\end{qoptions}

\qn{10} If $P\Rightarrow(q\vee r)$ is false, then the truth values of p,q,r are respectively:
\begin{qoptions}
\opt{A}{T,F,F}
\opt{B}{F,F,F}
\opt{C}{F,T,T}
\opt{D}{T,T,F}
\end{qoptions}

\qn{11} The value of\\ \resizebox{\linewidth}{!}{$\sin\frac\pi{14}\sin\frac{3\pi}{14}\sin\frac{5\pi}{14}\sin\frac{9\pi}{14}\sin\frac{11\pi}{14}\sin\frac{13\pi}{14}$}\\ is :
\begin{qoptions}
\opt{A}{$\dfrac1{16}$}
\opt{B}{$\dfrac1{64}$}
\opt{C}{$\dfrac1{128}$}
\opt{D}{None of these}
\end{qoptions}

\qn{12} Let $2\sin^2x+3\sin x-2>0$ and $x^2-x-2<0$ (x is in radians), then x lies in the interval $-$
\begin{qoptions}
\opt{A}{$\left(\dfrac\pi6,\dfrac{5\pi}6\right)$}
\opt{B}{$\left(-1,\dfrac{5\pi}6\right)$}
\opt{C}{$(-1,2)$}
\opt{D}{$\left(\dfrac\pi6,2\right)$}
\end{qoptions}

\qn{13} $\displaystyle\int_0^\pi\dfrac{x\tan x}{\sec x+\tan x}dx=$
\begin{qoptions}
\opt{A}{$\dfrac\pi2(\pi-1)$}
\opt{B}{$\dfrac\pi2\left(\dfrac\pi2-1\right)$}
\opt{C}{$\dfrac\pi2(\pi-2)$}
\opt{D}{$\dfrac\pi2\left(\dfrac\pi2-2\right)$}
\end{qoptions}

\qn{14} The area of the region\\ $\{(x,y):y^2\leq4x,\ 4x^2+4y^2\leq9\}$ is :-
\begin{qoptions}
\opt{A}{{\small $\dfrac{9\pi}8-\dfrac94\sin^{-1}\left(\dfrac13\right)+\dfrac{\sqrt2}6$}}
\opt{B}{{\small $\dfrac{9\pi}{16}-\dfrac98\sin^{-1}\left(\dfrac13\right)+\dfrac{\sqrt2}{12}$}}
\opt{C}{{\small $\dfrac{9\pi}4-\dfrac92\sin^{-1}\left(\dfrac13\right)+\dfrac{\sqrt2}3$}}
\opt{D}{{\small $\dfrac{9\pi}{32}-\dfrac9{16}\sin^{-1}\left(\dfrac13\right)+\dfrac{\sqrt2}{24}$}}
\end{qoptions}

\qn{15} A curve passes through the point $\left(1,\dfrac\pi6\right)$ Let the slope of the curve at each point (x, y) be $\dfrac yx+\sec\left(\dfrac yx\right)$, $x>0$. Then the equation of the curve is :-
\begin{qoptions}
\opt{A}{$\sin\left(\dfrac yx\right)=\log x+\dfrac12$}
\opt{B}{$\text{cosec}\left(\dfrac yx\right)=\log x+2$}
\opt{C}{$\sec\left(\dfrac{2y}x\right)=\log x+2$}
\opt{D}{$\cos\left(\dfrac{2y}x\right)=\log x+\dfrac12$}
\end{qoptions}

\qn{16} If $\hat a$ and $\hat b$ are unit vectors such that $\left[\hat a\ \hat b\ \hat a\times\hat b\right]=\dfrac14$, then angle between $\hat a$ and $\hat b$ is:-
\begin{qoptions}
\opt{A}{$\dfrac\pi3$}
\opt{B}{$\dfrac\pi4$}
\opt{C}{$\dfrac\pi6$}
\opt{D}{$\dfrac\pi2$}
\end{qoptions}

\qn{17} If $P\equiv(0,1,0)$, $Q\equiv(0,0,1)$, then projection of PQ on the plane $x+y+z=3$ is-
\begin{qoptions}
\opt{A}{$\sqrt3$}
\opt{B}{3}
\opt{C}{$\sqrt2$}
\opt{D}{2}
\end{qoptions}

\end{multicols}

\end{document}
